% SOURCE docs/release/2026-10-29/source/docs/writing-repair/2026-10-29/src/19bce971619f032a.tex CHAR 0
\chapter{Longitudinal interventions, state estimation and policy value}
\label{r3:longitudinal}
\section{One decision clock for the causal and statistical models}
A household's liquidity changes after an offer, its updated balance
influences the next offer, and both affect later funding contribution.
Collections has the same feedback: assistance changes arrears, arrears
affect subsequent assistance, and both determine repayment. The model
must preserve which information was available before each action.

Let decisions occur at $t=0,\ldots,T-1$. Before decision $t$, the bank
has observed $L_t$ and history
$H_t=(L_0,A_0,L_1,\ldots,A_{t-1},L_t)$. For now let $L_{t+1}$ include
the period reward $R_t$, updated measurements and eligibility, with
terminal outcome $Y$ contained in $L_T$. Behavior chooses an action with
probability $e_t(a\mid H_t)$. A target policy $g_t(a\mid H_t)$ can be
stochastic, or deterministic with all mass at $d_t(H_t)$. It must be
measurable with respect to the stated decision information and assign
only feasible actions. Define
\[
 \psi(g)=\E[Y^g],\qquad
 V(g)=\E^g\!\left[\sum_{t=0}^{T-1}\gamma^tR_t\right].
\]
The first is a terminal target; the second is an additive reward target.
The reward belongs to the period after the action, and the discount factor
converts later rewards to the initial decision date. Either target can be
encoded using an accumulated terminal reward, but the interpretation and
discount clock must remain explicit. The backward functions below initially
evaluate the terminal target; the later belief-state recursion evaluates
the additive reward target.
This chapter uses a finite action set. With continuous treatment,
probability sums become kernel integrals and weighting requires
Radon--Nikodym derivatives. A deterministic point policy is generally
singular relative to a continuous behavior density, so the weighting
and regular-inference formulas below do not carry over unchanged.

The observational chain rule always gives
\begin{equation}
 p(l_0,\bar a,\bar l_{1:T})
 =p(l_0)\prod_{t=0}^{T-1}e_t(a_t\mid h_t)
                              p(l_{t+1}\mid h_t,a_t).
 \label{r3:observed-sequence}
\end{equation}
No causality is asserted by this identity. A causal interpretation
requires consistency, no interference at the chosen unit, and sequential
exchangeability: conditional on each $H_t$, the action is independent of
the future counterfactual trajectory under the specified continuation
regimes, on compatible past histories. Additionally, every action/history
reached with positive probability under $g$ must have support under the
observed mechanism. This is a statement about sequential information,
not a claim that the sequence is first-order Markov.

\section{Mechanism replacement yields the g-formula}
\label{r3:g-formula}
Under these conditions, sequential exchangeability and consistency
identify the next-observation kernel under an intervened action with
the observed kernel conditional on its history. Induction from the
initial law therefore gives
\begin{equation}
 p^g(l_0,\bar a,\bar l_{1:T})
 =p(l_0)\prod_{t=0}^{T-1}g_t(a_t\mid h_t)
                              p(l_{t+1}\mid h_t,a_t).
 \label{r3:sequential-g}
\end{equation}
At the induction step, histories follow their already constructed
intervention law, the policy selects $a_t$, and the invariant next-data
kernel extends the history. This proves the factor replacement rather
than assuming it from a predictive fit. Summing $Y$ against this law
identifies $\psi(g)$ \citep[Chapter 21, Technical Point 21.1]{hernan2020}.

Equation~\eqref{r3:sequential-g} can be evaluated by simulation from
fitted next-data kernels. Alternatively define $V_T(H_T)=Y$ and recurse
backward:
\begin{align}
 Q_t(h,a)&=\E[V_{t+1}(H_{t+1})\mid H_t=h,A_t=a],\nonumber\\
 V_t(h)&=\sum_a g_t(a\mid h)Q_t(h,a),\qquad
 \psi(g)=\E[V_0(H_0)].
 \label{r3:sequential-recursion}
\end{align}
Iterated expectation under the intervention law proves the recursion.
This is a Bellman evaluation on the complete history, without a Markov
state reduction or an optimization claim. Sequential outcome regression
and forward g-computation evaluate the same functional using different
factorizations and approximation choices. They can have different
finite-sample behavior when the fitted models are restricted.

The change of measure between the two sequence laws is
\begin{equation}
 W_{T-1}^g=\prod_{t=0}^{T-1}
       \frac{g_t(A_t\mid H_t)}{e_t(A_t\mid H_t)},\qquad
 \psi(g)=\E[W_{T-1}^gY].
 \label{r3:sequence-weights}
\end{equation}
Multiplying Equation~\eqref{r3:observed-sequence} by the weight cancels
every observed assignment kernel and leaves
Equation~\eqref{r3:sequential-g}. The same cancellation proves
$\E[W_t^g]=1$ and the martingale property
$\E[W_t^g\mid H_t]=W_{t-1}^g$, taking $W_{-1}=1$.
Support is needed for the ratio; moments of the product are needed for
stable estimation. When a deterministic policy agrees with behavior
with probability $p$ at each of $T$ independent decisions, only a
fraction $p^T$ of trajectories follow it completely. Millions of rows
can coexist with little information for one sustained regime.

\section{Treatment--confounder feedback changes the regression target}
For an exact two-decision example, let
\begin{align*}
 P(L_0=1)&=0.5,\\
 P(A_0=1\mid L_0)&=0.2+0.6L_0,\\
 P(L_1=1\mid L_0,A_0)&=0.1+0.5L_0+0.2A_0,\\
 P(A_1=1\mid H_1)&=0.15+0.5L_1+0.1A_0,\\
 Y&=20+3A_0+4A_1+10L_1+6L_0.
\end{align*}
Independent assignment and outcome innovations implement these kernels.
Under sustained $A_0=A_1=1$, $\E[L_1]=0.55$ and $\E[Y]=35.5$.
Under $A_0=A_1=0$, they are $0.35$ and $26.5$, so the total effect is
$9$. Regressing $Y$ on all four inputs recovers the exact conditional
equation, but summing its action coefficients gives $7$. The missing
$2$ is the effect of $A_0$ mediated by $L_1$. Omitting $L_1$ instead
confounds the $A_1$ coefficient. Thus even a perfect conditional
regression requires the g-formula's averaging operation to answer the
sustained-intervention question.

The example has no hidden confounding, perfect measurement and support
for all action paths. Its failure is entirely a target error. This is
why the theoretical treatment must precede a banking pilot: the problem
can already be proved without institutional data.

\section{Marginal structural models and structural nested models}
\label{r3:msm-snm}
A marginal structural model (MSM) parameterizes a collection of regime
means, such as
$\E[Y^{a_0,a_1}]=\beta_0+\beta_1a_0+\beta_2a_1+\beta_{12}a_0a_1$.
An estimator must account for the assignment mechanism, rather than fit
that marginal mean directly to the observed conditional outcome law.
For fixed $\bar a$,
\[
 \E\!\left[\frac{\ind\{\bar A=\bar a\}}
                         {\prod_t e_t(a_t\mid H_t)}Y\right]
 =\E[Y^{\bar a}].
\]
Weighted estimating equations combine such identities across regimes.
If the MSM is misspecified, the result is a projection determined by
the chosen regime weights. Stabilized weights with numerators depending
on past actions create a different pseudopopulation; a correctly
specified MSM may retain its parameter, but a misspecified projection
and its variance can change. Weight clipping is another bias/variance
choice, not a neutral algebraic transformation
\citep[Chapter 21]{hernan2020}.

A structural nested mean model (SNMM) instead parameterizes the mean
effect of a treatment blip at a time, within the history and receipt
strata at that time, with subsequent treatment set to a reference regime.
For two decisions and reference zero, define
\begin{align}
 \gamma_1(h_1,a_1;\beta)
 &=\E[Y^{a_0,a_1}-Y^{a_0,0}\mid H_1=h_1,A_1=a_1],\nonumber\\
 \gamma_0(h_0,a_0;\beta)
 &=\E[Y^{a_0,0}-Y^{0,0}\mid H_0=h_0,A_0=a_0].
 \label{r3:blips}
\end{align}
The conditioning is on histories compatible with the displayed past;
these are mean contrasts and need not specify each person's effect.
At the true parameter let
$U_1=Y-\gamma_1(H_1,A_1;\beta)$ and
$U_0=U_1-\gamma_0(H_0,A_0;\beta)$.
By the definition of $\gamma_1$, the conditional mean of $U_1$ equals
that of $Y^{A_0,0}$. Sequential exchangeability makes that mean invariant
to $A_1$ given $H_1$. Averaging over the next history and removing
$\gamma_0$ then makes the mean of $U_0$ invariant to $A_0$ given $H_0$.
Consequently, for appropriate history instruments $d_t(H_t)$,
\begin{equation}
 \E\!\left[d_t(H_t)\{A_t-e_t(1\mid H_t)\}
           \{U_t(\beta)-\E[U_t(\beta)\mid H_t]\}\right]=0,
 \quad t=0,1.
 \label{r3:g-estimation}
\end{equation}
These moment equations define a g-estimation strategy. Identification
requires rank and a correct blip model. The mean restriction does not
assert that the blipped outcome equals the individual's no-treatment
counterfactual, nor that it is conditionally independent of treatment
in its entire distribution. Those are stronger distributional or
rank-preserving claims \citep[Sections 3.2, 5 and 6.1]{snm2014}.

In the exact example above, $\gamma_1=4a_1$ and $\gamma_0=5a_0$.
The early blip includes the $2$-unit mediated effect, unlike the
coefficient of $A_0$ in the regression conditioning on $L_1$.
These values make the sequential residual moments vanish. Replacing
the early blip by $3a_0$ leaves a nonzero moment. Thus g-estimation's
causal parameter is visible in a checkable moment calculation.

The relationship with IV is now precise. In Chapter~\ref{r3:instruments},
a valid instrument supplies a residual moment for an endogenous receipt.
Under sequential exchangeability, the residual of the action mechanism
supplies a moment after later effects have been removed. Structural
nested IV variants replace ignorability with other restrictions; their
identified blip parameters need not identify arbitrary marginal regime
means. An MSM, an SNMM and a predictive state-space model are different
models of different objects, even when all use linear functions.

\section{Recursive augmentation and what double robustness means}
\label{r3:long-dr}
For a fixed externally specified policy $g$, define $Q_t,V_t$ as in
Equation~\eqref{r3:sequential-recursion}. A mean-zero score is
\begin{equation}
 \phi_g=V_0(H_0)-\psi(g)+
      \sum_{t=0}^{T-1}W_t^g
       \{V_{t+1}(H_{t+1})-Q_t(H_t,A_t)\}.
 \label{r3:long-score}
\end{equation}
Each summand has zero expectation by conditioning on $(H_t,A_t)$.
This fact alone establishes an unbiased estimating expression at the
truth; it does not establish that the expression is a canonical gradient.
We next derive that stronger statement and its estimation remainder.

\subsection{A pathwise proof for an arbitrary fixed horizon}
\label{r3:long-gradient-proof}
Consider independent trajectories with a dominated, unrestricted
observed-data law. Assume the support needed for the target policy and
that the expressions below are square integrable. Along a regular
submodel with scalar index $\varepsilon$, the observed score decomposes as
\[
 S=S_0(L_0)+\sum_{t=0}^{T-1}
       \{S_{A,t}(A_t,H_t)+S_{L,t}(L_{t+1},H_t,A_t)\},
\]
where the action and next-observation scores have conditional means zero
given $H_t$ and $(H_t,A_t)$ respectively. These are mutually orthogonal
martingale increments in the observation order.

Keep $g$ fixed along the submodel. Differentiating the product in
Equation~\eqref{r3:sequential-g} gives
\[
 \dot\psi(g)=\E^g\!\left[Y\left\{S_0+\sum_t S_{L,t}\right\}\right].
\]
No action score appears because the intervention replaced that factor.
For the initial term, integrate the future under $g$ to obtain
$\E[(V_0-\psi)S_0]$. For a next-observation term, integrate the future
after $L_{t+1}$, then change measure only through action $t$:
\begin{align*}
 \E^g[YS_{L,t}]&=\E[W_t^gV_{t+1}S_{L,t}]\\
 &=\E[W_t^g(V_{t+1}-Q_t)S_{L,t}].
\end{align*}
The second equality uses the conditional mean-zero property of
$S_{L,t}$. Moreover $W_t^g(V_{t+1}-Q_t)$ itself lies in this
next-observation tangent component. Orthogonality of the components
therefore gives $\dot\psi(g)=\E[\phi_g S]$. In the unrestricted
model their sum spans $L_2^0(P)$, so Equation~\eqref{r3:long-score}
is the canonical gradient. This reconstructs the projection argument
for dynamic rules in \citet[Chapter 4, Section 4.4.1]{targeted2018},
here written for a fixed stochastic policy using the shared clock.
It also explains the nested-functional perspective of \citet{nested2022}.

The same orthogonality gives the efficiency-bound decomposition
\begin{equation}
 \Var(\phi_g)=\Var(V_0(H_0))+
 \sum_{t=0}^{T-1}\E\!\left[(W_t^g)^2
       \Var(V_{t+1}(H_{t+1})\mid H_t,A_t)\right].
 \label{r3:long-variance}
\end{equation}
For example, a policy that seldom agrees with historical collections
actions can inflate the late weighted innovation variance. Predicting
the initial risk accurately need not control that component. If the
displayed variance is infinite, the square-integrable gradient argument
does not deliver a regular root-$n$ estimate. Fixed finite horizon and
independent trajectories are part of the argument; neither an infinite
horizon nor common macroeconomic shocks are covered by changing notation.

It is useful to establish more than the true-nuisance identity. For any
fixed collection of candidate regressions $\widetilde Q_t$, define
$\widetilde V_t=\sum_a g_t\widetilde Q_t$, with
$\widetilde V_T=Y$. If all action kernels are correct, the expectation
of the augmented sum telescopes:
\[
 \E[\widetilde V_0]+\sum_t
  \{\E^g[\widetilde V_{t+1}]-\E^g[\widetilde V_t]\}
 =\E^g[Y].
\]
If all true $Q_t$ are used, every residual has zero conditional mean
even with other positive candidate action kernels. These are two valid
robustness statements. More refined sequential double robustness
requires the appropriate recursive construction; it cannot be assigned
to every list of separately fitted regressions
\citep{sequentialdr2018}.

The two-period remainder exposes the structure. Let
$q_t=\widetilde Q_t-Q_t$, $v_t=\widetilde V_t-V_t$,
$\widetilde W_0=g_0/\widetilde e_0$ and $v_2=0$.
Conditioning successively gives the exact bias
\begin{align}
 \E[\widetilde\Phi_g]-\psi(g)
 =&\ \E\!\left[\sum_a g_0(a\mid H_0)
        \left\{1-\frac{e_0(a\mid H_0)}{\widetilde e_0(a\mid H_0)}\right\}
                              q_0(H_0,a)\right]\nonumber\\
 &+\E\!\left[\widetilde W_0\sum_a g_1(a\mid H_1)
       \left\{1-\frac{e_1(a\mid H_1)}{\widetilde e_1(a\mid H_1)}\right\}
                              q_1(H_1,a)\right].
 \label{r3:two-period-remainder}
\end{align}
For example, the $v_1$ term in the first residual cancels the
$e_1/\widetilde e_1$ part of the second residual except for their
propensity-error product. The weights and history distributions in the
second term matter for rates. Longer horizons require their control as
well as nuisance accuracy. The sequential RL doubly robust estimator
is a closely related Bellman-residual construction, with its own
fixed-logging and independent-trajectory assumptions
\citep[Section 4.1]{jiang2016}.

\subsection{General remainder and the precise robustness condition}
\label{r3:long-remainder-proof}
Define $q_t=\widetilde Q_t-Q_t$, $v_t=\sum_a g_tq_t$,
$v_T=0$ and $\widetilde W_t=\prod_{s=0}^tg_s/\widetilde e_s$,
with $\widetilde W_{-1}=1$. Candidate functions are fixed here;
for fitted functions the identity holds conditionally on independent
training data. Subtract the true augmented expression. Each true
residual has conditional mean zero even after multiplication by the
candidate history/action weight. The remaining bias is
\[
 B=\E[v_0]+\sum_{t=0}^{T-1}
       \E[\widetilde W_t(v_{t+1}-q_t)].
\]
Relabel the positive terms using $v_T=0$ and then integrate the current
action in the negative terms. This proves the general identity
\begin{equation}
 B=\sum_{t=0}^{T-1}\E\!\left[
 \widetilde W_{t-1}\sum_a g_t(a\mid H_t)
 \frac{\widetilde e_t(a\mid H_t)-e_t(a\mid H_t)}
      {\widetilde e_t(a\mid H_t)}q_t(H_t,a)\right].
 \label{r3:long-remainder}
\end{equation}
This is more informative than a statement that ``the estimator is DR.''
Each product vanishes if that time's propensity is correct or that
time's candidate regression is the true continuation $Q_t$. However,
the true continuation already incorporates every later target action.
A correctly specified regression of a biased downstream pseudo-outcome
does not necessarily estimate it. Thus an alternating pattern of
correctly specified regression models and propensity models is not
automatically sufficient for an arbitrary fitting procedure.

One construction uses backward augmented pseudo-outcomes:
\begin{equation}
 D_T=Y,\qquad
 D_t=\widetilde V_t(H_t)+
       \frac{g_t(A_t\mid H_t)}{\widetilde e_t(A_t\mid H_t)}
       \{D_{t+1}-\widetilde Q_t(H_t,A_t)\}.
 \label{r3:sdr-recursion}
\end{equation}
Suppose by induction that
$\E[D_{t+1}\mid H_t,A_t]=Q_t(H_t,A_t)$, as follows when the
next history conditional mean is the true $V_{t+1}$.
Then $\E[D_t\mid H_t]=V_t(H_t)$ if either
$\widetilde e_t=e_t$ or $\widetilde Q_t=Q_t$.
The base case is the terminal observed outcome. Regressing the
appropriate augmented continuation rather than an uncorrected fitted
future response can therefore support the backward induction, given
consistent conditional-mean learning and support. Expanding $D_0$
recovers the augmented sum. This is the unbiased-transformation route
discussed in \citet[Definitions 1--3 and Section 3]{sequentialdr2018};
that paper also develops a different, bounded targeting construction.
Pseudo-outcomes $D_t$ can lie outside the bounds on $Y$, especially
under small propensities, so they are not interchangeable with bounded
outcomes in a logistic regression package.

For a rate statement, let $\nu_t$ denote the law with the observed
$H_t$ marginal and action drawn from $g_t(\cdot\mid H_t)$. If
$\widetilde e_t\ge\epsilon$ on target actions, then
$\widetilde W_{t-1}\le\epsilon^{-t}$ and Cauchy--Schwarz gives
\begin{equation}
 |B|\le\sum_{t=0}^{T-1}\epsilon^{-(t+1)}
 \|\widetilde e_t-e_t\|_{L_2(\nu_t)}
 \|\widetilde Q_t-Q_t\|_{L_2(\nu_t)}.
 \label{r3:long-rate-bound}
\end{equation}
For fixed $T$, a vanishing $o_p(n^{-1/2})$ bound, convergence of the
whole fitted score in $L_2(P)$, and an appropriate independent-trajectory
CLT yield the asymptotically linear expansion for a cross-fitted
augmented estimator. The constants display why a growing-horizon
theorem cannot be obtained by reusing a fixed-horizon nuisance rate.
Truncating weights changes the estimator's bias; it does not change the
underlying intervention's support or prove the remainder negligible.

\subsection{Backward targeting is a sequence of score equations}
\label{r3:long-targeting}
For a concrete bounded-outcome construction, suppose $0\le Y\le1$,
$g$ is fixed, and candidate propensities are positive. Begin at
$V_T^*=Y$ and proceed backward. At time $t$, regress the updated
$V_{t+1}^*$ on $(H_t,A_t)$ to obtain $\widehat Q_t$ in $(0,1)$.
Hold the later update fixed and minimize Bernoulli cross-entropy with
fractional response $V_{t+1}^*$ along
\begin{equation}
 \operatorname{logit}Q_{t,\varepsilon}(h,a)
 =\operatorname{logit}\widehat Q_t(h,a)
       +\varepsilon C_t(h,a),\qquad
 C_t(h,a)=\widehat W_{t-1}(h)
                   \frac{g_t(a\mid h)}{\widehat e_t(a\mid h)}.
 \label{r3:long-fluctuation}
\end{equation}
At a finite minimizer its likelihood score is
$P_n[C_t(V_{t+1}^*-Q_t^*)]=0$. Set
$V_t^*(h)=\sum_a g_t(a\mid h)Q_t^*(h,a)$ and use it as the
response at the preceding step. Finally let
$\widehat\psi=P_n V_0^*$. Later equations stay solved because a
new earlier update changes neither their fitted continuation nor their
weights. Summing the time-specific score equations therefore gives
$P_n\widehat\phi_g=0$. The estimate remains in $[0,1]$.
This is the same fluctuation principle as the dynamic-rule algorithm
of \citet[Chapter 4, Sections 4.4.2--4.4.3]{targeted2018}, expressed
here with stochastic policy averages and full histories.

A finite optimizer, empirical score tolerance, compatible rule/history
evaluation and specified training are necessary implementation details.
Perfect separation can send a logistic fluctuation to infinity; a
numerical early stop is then not evidence that the score is solved.
Efficiency additionally requires the updated nuisance remainder and
empirical-process conditions, as in
\citet[Chapter 4, Section 4.5]{targeted2018}. The bounded targeting
construction is not, by its name alone, the stronger sequentially
robust fitting algorithm above. The product must record the actual
recursion and guarantee used.

The derivations concern a fixed policy. If the intervention itself
depends on the observed propensity, for example
$g_{\delta,t}(1\mid h)=\delta e_t(1\mid h)/\{1-e_t(1\mid h)+\delta e_t(1\mid h)\}$,
then perturbing $P$ also perturbs $g_{\delta,t}$. Its gradient has
additional terms \citep{kennedy2019}. Treating that policy as fixed
would omit part of the target's derivative.

\section{Observation, censoring and eligibility are mechanisms}
Suppose $C_{t+1}=1$ means the next observation is retained, with
$\pi_t=P(C_{t+1}=1\mid H_t,A_t,C_t=1)$ and $C_0=1$. Under the
appropriate conditional independent-censoring assumption and support,
the weight for a terminal observed outcome additionally includes
\[
 \prod_{t=0}^{T-1}\frac{C_{t+1}}{\pi_t(H_t,A_t)}.
\]
This is a second mechanism replacement: it emulates retaining the
defined outcome. It does not emulate preventing default, death or
another event that changes what the endpoint means. Re-entry, debt sale,
account closure and competing events need separate clocks and targets.
Eligibility $E_t$ affected by prior actions belongs in $H_t$ and in
the feasible action set. Restricting retrospectively to people eligible
at every time selects on post-treatment history. A policy must specify
what happens on ineligible histories rather than delete them silently.

\section{A controlled state-space model on the same clock}
\label{r3:state-bridge}
Now suppose the relevant condition is a hidden state $S_t$ and the bank
observes a signal $O_t$ before acting. Let
$H_t=(O_0,A_0,R_0,O_1,\ldots,A_{t-1},R_{t-1},O_t)$.
A particular controlled model supplies initial law $\mu_\theta$,
initial emission $B_{0,\theta}$, transition
$P_\theta(s'\mid s,a)$, reward density
$F_\theta(r\mid s,a,s')$, and next signal density
$B_\theta(o'\mid s',a)$. This last conditional-independence structure
is a model restriction. A joint reward/signal density is needed if
those emissions are not conditionally independent as specified.

Let $q_\eta(a\mid s,h)$ represent behavior, possibly using private
information summarized by $s$. With other time-varying inputs understood
as explicit arguments, the complete trajectory law is
\begin{align}
 p_{\theta,\eta}(s_{0:T},o_{0:T},a_{0:T-1},r_{0:T-1})
 =&\ \mu_\theta(s_0)B_{0,\theta}(o_0\mid s_0)\nonumber\\
 &\times\prod_{t=0}^{T-1}q_\eta(a_t\mid s_t,h_t)
       P_\theta(s_{t+1}\mid s_t,a_t)\nonumber\\
 &\hspace{13mm}\cdot F_\theta(r_t\mid s_t,a_t,s_{t+1})
       B_\theta(o_{t+1}\mid s_{t+1},a_t).
 \label{r3:controlled-law}
\end{align}
The observed likelihood sums this law over the hidden states. A causal
interpretation further requires the state-transition, reward and signal
mechanisms to remain invariant when the action mechanism is replaced,
with state sufficient to capture relevant hidden dependence. Those
restrictions can be expressed by a time-indexed SCM with mutually
independent innovations conditional on the declared state. A model
fitted to the observational likelihood does not prove that it is such
an SCM.

Under a target policy using only recorded history, replace the $q_\eta$
factor by $g_t(a\mid h)$. This defines a model-conditional intervention
law. Whether its value is identified from historical observations is a
separate question about the model class and information, to which we
return below. A private logger with memory not captured by $S_t,H_t$
requires state augmentation or a more general behavior kernel; the
displayed finite-state model would otherwise be misspecified.

\subsection{Recovering the observed-history g-formula}
If $q_\eta(a\mid s,h)=e_t(a\mid h)$ and the structural innovations
justify sequential exchangeability, actions supply no extra information
about $S_t$ conditional on $H_t$. Summing
Equation~\eqref{r3:controlled-law} over states defines observable history
kernels for $(R_t,O_{t+1})$. Under common supported histories, their
conditional state mixtures are invariant to replacing the observed-history
policy, because the action factors cancel from the latent posterior.
Hence replacing $e_t$ by $g_t$ yields the history g-formula.

This establishes the bridge: a controlled state-space model is a
restricted parameterization of longitudinal observation kernels when
its action and observation assumptions permit that reduction. It can
be computationally valuable when histories are long. It is not a weaker
replacement for the causal assumptions underlying those kernels.
When $S_t$ is observed and Markov sufficient, $Q_t(H_t,a)$ reduces to
$Q_t(S_t,a)$; under partial observation it can reduce to a belief state
under further restrictions.

\section{Filtering actions is different from intervening on them}
\label{r3:filtering}
Let $b_t(s)=P(S_t=s\mid H_t)$ before action. Under private logging,
observing action $a$ first updates the state distribution to
\begin{equation}
 b_t^{a,\mathrm{obs}}(s)
   =\frac{q_\eta(a\mid s,H_t)b_t(s)}
          {\sum_u q_\eta(a\mid u,H_t)b_t(u)}.
 \label{r3:action-update}
\end{equation}
After reward $r$ and signal $o'$, filtering gives
\begin{equation}
 b_{t+1}(s')=
 \frac{B_\theta(o'\mid s',a)
       \sum_s b_t^{a,\mathrm{obs}}(s)P_\theta(s'\mid s,a)
                                    F_\theta(r\mid s,a,s')}
      {\sum_{u,s}B_\theta(o'\mid u,a)b_t^{a,\mathrm{obs}}(s)
            P_\theta(u\mid s,a)F_\theta(r\mid s,a,u)}.
 \label{r3:controlled-filter}
\end{equation}
Both equalities are Bayes' rule applied to the joint law. Under
$\doop(A_t=a)$ there is no action-likelihood update: use $b_t$ in
place of $b_t^{a,\mathrm{obs}}$. With observed-history logging the two
beliefs coincide because $q_\eta$ is constant across $s$.

For an exact example, let $S$ be binary with prior probability $0.5$,
$P(A=1\mid S=1)=0.8$, $P(A=1\mid S=0)=0.2$, and
$P(Y=1\mid S,A)=0.2+0.6S$. Then $P(S=1\mid A=1)=0.8$ and
$P(S=1\mid A=0)=0.2$, giving observed outcome means $0.68$ and
$0.32$. Intervening on either action retains the state prior and gives
mean $0.5$, so the true effect is zero. A correct observational filter
must use the action as evidence; a correct intervention calculation
must remove that selection factor. Filtering the hidden condition more
accurately does not make the two operations equivalent.

Let
$K_t(s,s')=q_\eta(a_t\mid s,h_t)P_\theta(s'\mid s,a_t)
F_\theta(r_t\mid s,a_t,s')B_\theta(o_{t+1}\mid s',a_t)$.
Normalized forward messages satisfy
\[
 \widetilde\alpha_{t+1}(s')=\sum_s\alpha_t(s)K_t(s,s'),\quad
 c_{t+1}=\sum_{s'}\widetilde\alpha_{t+1}(s'),\quad
 \alpha_{t+1}=\widetilde\alpha_{t+1}/c_{t+1}.
\]
Starting from $\alpha_0(s)\propto\mu(s)B_0(o_0\mid s)$, the
log likelihood is $\sum_{t=0}^T\log c_t$. The initial normalizing
constant is the probability of the first signal under the initial state
mixture. Backward messages and smoothed transition counts yield the usual
EM construction, now for the declared controlled kernel \citep{rabiner1989}.

If $q_\eta=e_\eta(a\mid h)$ and its parameters are distinct from
$\theta$, its factors can be removed from the objective for estimating
$\theta$ because they factor outside the latent-state sum. This is
a factorization of the joint likelihood. It does not say that the
remaining factor equals the conditional density of all observations
given the entire adaptively chosen action sequence. Future actions can
depend on earlier observations. Conditioning on their whole sequence
would require a different normalization.

Smoothing estimates $P(S_t\mid H_T)$ and uses later information.
It is useful for parameter estimation and historical diagnosis, but
is not a state available to the decision at $t$. Using it as if it
were a prospective policy input creates leakage. When parameters are
uncertain, a Bayesian sufficient state may be the joint posterior of
$(S_t,\theta)$; plugging in $\widehat\theta$ produces a conditional
planning calculation and omits parameter-learning value and uncertainty.

\section{Belief-state planning and three kinds of identification}
\label{r3:belief-value}
Assume a known, invariant controlled Markov environment and a target
policy using the full filtering belief, with all other relevant observed
inputs included. Define
\[
 r(b,a)=\sum_{s,s'}b(s)P_\theta(s'\mid s,a)
                              \E_\theta[R\mid s,a,s'].
\]
The distribution of the next reward/signal depends on history only
through $b,a$, and the next belief is a deterministic Bayes update
$\mathcal B(b,a,r,o')$. The value recursion is therefore
\begin{equation}
 V_t^g(b)=\sum_a g_t(a\mid b)
       \left[r(b,a)+\gamma\E_{r,o'\mid b,\doop(a)}
            V_{t+1}^g(\mathcal B(b,a,r,o'))\right],
 \quad V_T^g=0.
 \label{r3:belief-bellman}
\end{equation}
Replacing the policy average by a maximum gives the finite-horizon
planning recursion when actions are feasible and rewards integrable.
For a stationary discounted model with bounded rewards and
$0<\gamma<1$, the optimal Bellman operator is a contraction in sup norm:
the difference of two maximized continuation values is at most
$\gamma\|V-W\|_\infty$. This gives a unique bounded fixed point.
The belief-state sufficiency construction is developed in
\citet[Sections 3.2--3.4]{kaelbling1998}; it assumes a supplied model.

Three distinct identification questions precede causal use of the value.
First, can the observed law identify the likelihood parameters within
the proposed state-space class? For an uncontrolled stationary $K$-state
HMM, full-rank ergodic transition and distinct emission distributions
identify parameters up to labels from a sufficiently long observation
block under the cited conditions
\citep[Theorem 1]{hmmident2016}. This theorem is not automatically a
controlled, privately logged HMM theorem.

Second, does the historical model identify the intervention law after
changing $q$? Causal autonomy, action support and restrictions on hidden
assignment and measurement are additional. Third, is the requested
policy value uniquely determined even if some parameters are not?
Formally, for every two admissible models with the same observed law,
one needs $V_{\mathcal M_1}(g)=V_{\mathcal M_2}(g)$, which is weaker
than equality of every latent parameter but stronger than good fit.

A permutation of hidden-state labels leaves both observation likelihood
and coherently computed policy values unchanged; it is a representational
ambiguity. The two worlds in Chapter~\ref{r3:graphs} instead have the
same likelihood and different action values; that is causal ambiguity.
A bootstrap conditional on one selected model measures sampling
variation within it and does not cover the latter ambiguity. Specialized
POMDP off-policy identification uses informative proxies or other
restrictions \citep{tennenholtz2020}; it is not supplied by a generic
filter or by the contraction proof.

\section{What a state-space adapter must preserve}
\label{r3:state-contract}
A numerical filter consumes transition and observation kernels, initial
state, parameter values and masks. A causal adapter additionally needs
the action clock, logger information, intervention mechanism, reward
definition and source of identification. It must specify whether a
missing measurement is exogenous, conditionally ignorable, or modeled
through a state-dependent observation process. Omitting an endogenous
mask changes the likelihood and may change policy identification.

For continuous states, sums become integrals and Kalman, nonlinear
Gaussian, particle or variational methods approximate different objects.
The state-space adapter should report the actual likelihood approximation,
not only the latent-state dimension. Label fixing, deterministic lag
coordinates, parameter priors and reward integration must be coherent
across fitting and intervention evaluation. This is how a numerical
filtering library can be reused without attributing causal guarantees
to its numerical accuracy.

The customer and the bank may also be different decision makers. The
bank's POMDP maximizes a supplied bank reward. A model of customer
activation, refinancing or default requires utility, beliefs, information
and transition restrictions of its own. Estimating a customer's dynamic
choice likelihood is not equivalent to optimizing the bank's reward;
equilibrium response adds another layer. Chapter~\ref{r3:controlled-identification}
develops that inverse decision problem; Chapters~\ref{r3:time-series}
and \ref{r3:continuous} connect structural shocks and feedback to it.

